\documentclass[aspectratio=169,dvipsnames,table]{beamer}

% Basic packages
\usepackage{hyperref}

% Absolute position for text and elements
\usepackage[absolute,overlay]{textpos}

% Use latim special characters
\usepackage[utf8]{inputenc}

% DejaVu Sans (Type 1-converted, works under plain pdflatex -- no need for
% xelatex/lualatex) -- the same font the companion software
% (software/efficient_nn_lab) already renders in its matplotlib panels and
% Qt UI, so slides and live software read as one visual system. Sets
% \sfdefault; beamer already routes body text through \sfdefault by
% default, so no \familydefault override is needed.
\usepackage{DejaVuSans}

% Use colors
\usepackage{color}

% Enables graphics
\usepackage{graphicx}

% Better text justification
\usepackage{microtype}

% Enables multicolumn text
\usepackage{multicol}

% Enables \toprule/\midrule/\bottomrule in tables
\usepackage{booktabs}

% Fine-grained control over list spacing (dense two-column slides)
\usepackage{enumitem}
\setlist{itemsep=2pt, topsep=2pt, parsep=0pt}

% Enables math
\usepackage{amsthm}
\usepackage{amsmath}
\usepackage{amssymb}

% Adds some symbols (like slash on a non mathematical context)
\usepackage{textcomp}

% Quote packages, ABNT standard
\usepackage[alf]{abntex2cite}

% Portuguese-specific commands
\usepackage[brazilian]{babel}

% Bibliography
\usepackage[brazilian,hyperpageref]{backref}

% Drawing and plot packages
\usepackage{relsize}
\usepackage{tikz}
\usetikzlibrary{calc}
\usetikzlibrary{positioning}
\usetikzlibrary{decorations.pathreplacing}
\usetikzlibrary{shapes}
\usetikzlibrary{backgrounds}
\usetikzlibrary{patterns}
\usetikzlibrary{fit}
\usetikzlibrary{arrows.meta}

\usepackage{pgfplots}
\usepackage{pgfplotstable}
\pgfplotsset{compat=newest}

% Reuse the thesis' CVD-safe Okabe-Ito palette and diagram semantic roles
% (colors + \tableHeaderRow/\tableZebraStripe) instead of redefining our own.
% visualStyle.tex — Single source of truth for every color used in the
% thesis: tables, the Gantt chart, and TikZ diagrams (images/*.tex).
%
% Palette: Okabe & Ito (2008), "Color Universal Design" — the standard
% categorical palette distinguishable under protanopia, deuteranopia and
% tritanopia. Every color in the thesis must come from here; do not
% \definecolor or use raw color names (fill=orange, fill=green, fill=red, ...)
% ad hoc inside chapters/, tables/*.tex or images/*.tex. Raw xcolor names in
% particular (plain "red"/"green"/"yellow") are NOT CVD-safe — red/green is
% the single most common color-blindness confusion pair.

% ── Okabe-Ito base palette ────────────────────────────────────────────────
\definecolor{cbOrange}{RGB}{230,159,0}
\definecolor{cbSkyBlue}{RGB}{86,180,233}
\definecolor{cbGreen}{RGB}{0,158,115}     % bluish green
\definecolor{cbYellow}{RGB}{240,228,66}
\definecolor{cbBlue}{RGB}{0,114,178}
\definecolor{cbVermillion}{RGB}{213,94,0}
\definecolor{cbPurple}{RGB}{204,121,167}
\definecolor{cbGrey}{RGB}{166,166,166}

% ── Semantic roles (used by every table in the thesis) ───────────────────
% Tint levels below are chosen against two constraints, not eyeballed:
%  1. WCAG 1.4.11 (non-text contrast, >=3:1) / 1.4.3 (text contrast, >=4.5:1)
%     — verified against black text for the band colors below (~9-13:1,
%     comfortably AAA) and against white text for the solid header (~5.2:1,
%     AA). A pale tint (e.g. !35 or lower) still clears contrast but reads
%     as washed-out/grey on paper — the "too dull" failure mode — because
%     contrast alone doesn't capture perceived saturation/vividness.
%  2. Okabe & Ito (2008) hue choices are kept exactly (blue/orange is the
%     most CVD-robust 2-hue pair); only the tint (mix-with-white) percentage
%     is raised for a visibly saturated, "colorful but not extravagant" look.
%
% Header row: solid (untinted) blue background, white bold text.
\colorlet{tableHeaderColor}{cbBlue}
% Zebra striping on long/data-heavy tables: a visible-but-subtle tint of the
% header hue — enough to guide the eye across a row without competing with
% the semantic band colors below. Kept lighter than those on purpose: zebra
% striping is a navigation aid, not information-bearing.
\colorlet{tableAltRowColor}{cbBlue!12}
% Wavelet coefficient examples (tables/{regular,packet}WaveletExample.tex):
% grey = raw signal label, sky blue = row/column index label,
% blue = approximation band (cA), orange = detail band (cD).
% Blue/orange is the single most robust 2-hue pair across all CVD types.
% These ARE information-bearing (which band a cell belongs to), so they are
% tinted much more strongly than the zebra stripe above.
\colorlet{tableSignalLabelColor}{cbGrey}
\colorlet{tableIndexLabelColor}{cbSkyBlue}
\colorlet{tableSignalValuesColor}{cbGreen!55}
\colorlet{tableApproxBandColor}{cbBlue!55}
\colorlet{tableDetailBandColor}{cbOrange!55}

% Gantt chart (chapters/08-proposedApproach.tex, tab:cronogramaGantt):
% same palette, kept here so the chart matches the tables it sits beside.
\colorlet{ganttBarColor}{cbGreen}
\colorlet{ganttMilestoneColor}{cbGreen}
\colorlet{ganttTitleColor}{cbBlue!45}
\colorlet{ganttHighlightColor}{cbOrange}

% ── Reusable header row ────────────────────────────────────────────────────
% Usage inside a tabular/longtable, as the first thing in a header row:
%   \tableHeaderRow \textbf{A} & \textbf{B} \\
\newcommand{\tableHeaderRow}{\rowcolor{tableHeaderColor}\color{white}}

% ── Zebra striping for table/longtable bodies ─────────────────────────────
% Call once, right before \begin{tabular}/\begin{longtable}. Colors rows
% 2, 4, 6, ... (i.e. every data row after a 1-row header) with
% tableAltRowColor; row 1 (the header) is left untouched so \tableHeaderRow
% still controls it. NOTE: \rowcolor itself must always be the literal first
% token of a row — never wrap it in another macro (e.g. a counter-stepper),
% since colortbl's \noalign trick requires nothing to precede it.
\newcommand{\tableZebraStripe}{\rowcolors{2}{white}{tableAltRowColor}}

% ── Wide/wrapping text column, vertically centered ────────────────────────
% Plain p{<width>} columns top-align their content, while ordinary c/l/r
% columns baseline/vertically-center theirs — invisible in a plain \hline
% table, but a visible split (header text sitting below the colored fill)
% once a \rowcolor-filled header row mixes p{} with c/l/r columns. Use the
% column type Y{<width>} instead of p{<width>} in any column spec that will
% share a row with \tableHeaderRow, e.g. \begin{tabular}{@{}Y{10cm}cc@{}}.
\newcolumntype{Y}[1]{>{\raggedright\arraybackslash}m{#1}}

% ── Diagram / flowchart semantic roles (images/*.tex) ─────────────────────
% The thesis has ~20 TikZ flowcharts and node diagrams that all reuse the
% same handful of roles (start/end bubble, process box, decision diamond,
% arrows, autoencoder layer nodes). They used to \definecolor / fill= with
% raw xcolor names (orange, green, yellow, red, gray) independently per
% file — not CVD-safe (plain red/green/yellow is the classic confusion
% triple) and impossible to re-theme consistently. Every images/*.tex file
% must reference these instead of a raw color name.
\colorlet{flowStartEndColor}{cbOrange}
\colorlet{flowProcessColor}{cbGreen}
\colorlet{flowDecisionColor}{cbYellow}
\colorlet{flowArrowColor}{cbGrey}
% Autoencoder-style node diagrams (images/{autoencoder,denoisingAutoencoder,residuo}*.tex):
% grey = input, yellow = hidden/code/latent layer, vermillion = output/reconstruction.
% (Vermillion replaces the original "red" — true red is excluded from the
% Okabe-Ito set for the same red/green confusion reason.)
\colorlet{diagramInputColor}{cbGrey}
\colorlet{diagramHiddenColor}{cbYellow}
\colorlet{diagramOutputColor}{cbVermillion}
% Intermediate/corrupted stage between input and hidden layer (e.g. the
% noised x_hat in images/denoisingAutoencoder.tex, or the pre-activation
% sum in images/residuo.tex) — kept visually distinct from all three roles
% above.
\colorlet{diagramIntermediateColor}{cbSkyBlue}
% Tree diagrams (images/afasias.tex): darker tints so white text keeps
% AA contrast on non-header-sized diagram text.
\colorlet{diagramBranchColor}{cbGreen!80!black}
\colorlet{diagramLeafColor}{cbBlue!80!black}
% Feature-vector pipeline diagrams (images/{bark,mel,linear,lfcc,bfcc}...Vect.tex,
% images/{mel,linear}BandEnergies.tex): single accent color for box/ellipse
% outlines. Okabe-Ito blue, darkened for outline contrast against white fill.
\colorlet{diagramOutlineColor}{cbBlue!80!black}

% ── Marcação temporária de revisão (REMOVER antes da versão final) ────────
% Destaca um trecho que ainda depende de trabalho pendente, para que não
% passe despercebido em uma releitura. Usa cbYellow da paleta Okabe-Ito (e
% não o "yellow" bruto do xcolor, que não é CVD-safe); o tom é rebaixado
% para !40 porque o texto impresso por cima é preto e precisa manter
% contraste AA.
%
% \pendenteMarca{...} destaca um trecho CURTO e sem matemática: usa
% \colorbox, que não quebra linha nem sobrevive a $...$ — por isso não deve
% envolver parágrafos inteiros nem fórmulas. Para blocos, use \pendenteNota,
% que colore apenas o rótulo e deixa o corpo do texto em itálico comum.
\colorlet{pendenteColor}{cbYellow!40}
\newcommand{\pendenteMarca}[1]{\colorbox{pendenteColor}{#1}}
\newcommand{\pendenteNota}[1]{%
	\par\noindent\colorbox{pendenteColor}{\textbf{[PENDENTE]}}\hspace{0.4em}\textit{#1}\par}


% Semantic roles for this talk's two "frontiers" -- exact hex values from
% software/efficient_nn_lab/src/efficient_nn_lab/app/theme.py (BITNET_COLOR,
% SNN_COLOR, CONVERGE_COLOR, NEUTRAL_COLOR, ACCENT_COLOR), deliberately
% brighter/more saturated than the thesis' visualStyle.tex tints (those are
% calibrated for AA/AAA contrast in printed text; this talk is projected,
% not printed, and needs to match the live software instead). Do not
% propagate these values back into visualStyle.tex or the thesis chapters.
% bitnetColor/snnColor are darkened just enough (WCAG contrast >=4.5:1 on
% white) to stay legible where motivacao.tex/conclusao.tex use them as
% raw inline \textcolor{...}{} text on the plain white slide background --
% see theme.py's matching comment for the contrast math (same values,
% mirrored). convergeColor/accentColor are never used as raw inline text
% in this deck (only as light !NN!white box-background tints, always
% paired with dark chromeText), so they keep the fully vivid values.
\definecolor{bitnetColor}{HTML}{0073E5}
\definecolor{snnColor}{HTML}{E61F00}
\definecolor{convergeColor}{HTML}{00E38A}
\definecolor{neutralColor}{HTML}{5B6472}
\definecolor{accentColor}{HTML}{FFC300}

% Matches TEXT_COLOR in software/efficient_nn_lab/.../app/theme.py. Used
% as the fixed text color on every colored chrome bar/box below instead
% of white or a hue-matched color: measured contrast ratios (WCAG
% relative-luminance formula) against all five tinted title backgrounds
% range 6.1:1-14.5:1 with near-black text, comfortably clearing the 4.5:1
% AA threshold everywhere. White text only clears >=3:1 (AA-large) on the
% two darkest tints (bitnet/snn) and fails outright on converge/accent/
% neutral; a hue-matched color (e.g. bitnetColor text on a
% bitnetColor!92!white box) is even worse, since text and background
% converge toward the same color -- that was the previous, broken
% behavior in taxonomia.tex/desafios.tex's comparison boxes.
\definecolor{chromeText}{HTML}{111318}

% Packages configuration

\setbeamercolor{background canvas}{bg=}
\setbeamercolor{normal text}{bg=black!20}

% beamer's default "structure" color (a fixed blue) otherwise tints every
% itemize bullet, block title and section number in the deck blue,
% regardless of which of the two frontiers (BitNet=blue, SNN=red,
% convergência=green) is on screen. Pin it neutral so only the frametitle
% bar (set per-section below, in apresentacao.tex, right before each
% \section) and explicit \textcolor{...Color}{} calls carry a hue --
% that is what actually varies the palette across the deck instead of
% bathing everything in blue.
\setbeamercolor{structure}{fg=neutralColor}
\setbeamercolor{title}{fg=chromeText,bg=bitnetColor!92!white}
\setbeamercolor{frametitle}{fg=chromeText,bg=bitnetColor!92!white}

% Dedicated (non-dynamic) title-box styles for frames that show two
% frontiers side by side (taxonomia.tex, desafios.tex) -- these must NOT
% follow the per-section frametitle color below, since both colors need
% to be on screen at once.
\setbeamercolor{titleBitnet}{fg=chromeText,bg=bitnetColor!92!white}
\setbeamercolor{titleSnn}{fg=chromeText,bg=snnColor!85!white}
\setbeamercolor{titleConverge}{fg=chromeText,bg=convergeColor!80!white}
\setbeamercolor{titleAccent}{fg=chromeText,bg=accentColor!75!white}

% Clickable "open this demo in the software" cue -- presentation.md phase
% 3/4. #1 = link text (demo name as it appears in the software's sidebar),
% #2 = DemoModule.slug (see software/efficient_nn_lab/src/efficient_nn_lab/
% core/demo.py). "run:" paths resolve against this PDF's own directory, so
% abrir-demo.sh (next to this preamble) does the rest. Many PDF viewers
% block "run:" links by default (see presentation.md 4.2) -- the visible
% text is deliberately also the exact demo name, so it still works as a
% pure navigation label ("go click this in the sidebar") even where the
% link itself does nothing.
\newcommand{\abrirNoSoftware}[2]{%
	\par\vspace{2pt}\noindent%
	{\scriptsize\color{neutralColor}%
		$\blacktriangleright$~\href{run:./abrir-demo.sh #2}{Abrir no software: \textbf{#1}}%
	}%
}

% Enables section slides
\AtBeginSection[]{
	\begin{frame}
		\vfill
		\centering
		\begin{beamercolorbox}[sep=18pt,center,shadow=true,rounded=true]{title}
			\usebeamerfont{title}\insertsectionhead\par
		\end{beamercolorbox}
		\vfill
	\end{frame}
}


\title{As fronteiras da arquitetura de redes neurais}
\subtitle{BitNets e redes de pulso}

\institute{Universidade Estadual Paulista Júlio de Mesquita Filho}

\author{André Furlan - ensismoebius@gmail.com}

\logo{
\includegraphics[height=1cm]{unesp.png}}
\date{\the\year}

\begin{document}

	\frame{\titlepage}

	% Frametitle/title bar color changes per section, so the chrome itself
	% signals which frontier is on screen instead of staying blue
	% throughout (see preamble.tex's \setbeamercolor{structure} comment).
	% Must be set BEFORE each \section command: \AtBeginSection's divider
	% frame is emitted during \section's own expansion, so setting the
	% color only after \section would leave the divider frame one step
	% behind.
	\setbeamercolor{title}{fg=chromeText,bg=neutralColor!55!white}
	\setbeamercolor{frametitle}{fg=chromeText,bg=neutralColor!55!white}
	\section{Fundamentos}
		\begin{frame}[shrink]
	\frametitle{Como uma rede neural é construída}
	\begin{itemize}
		\item Uma rede neural é organizada em \textbf{camadas}: entrada $\rightarrow$ uma ou mais camadas ocultas $\rightarrow$ saída.
		\item Cada neurônio de uma camada recebe \textbf{todas} as saídas da camada anterior, cada uma multiplicada pelo seu próprio \textbf{peso}, soma tudo, e passa o resultado por uma \textbf{função de ativação} não linear.
		\item Sem a não linearidade, empilhar camadas equivaleria a uma única camada linear --- é a ativação que dá à rede o poder de aproximar funções complexas.
	\end{itemize}
	\vspace{4pt}
	\begin{center}
		\begin{tikzpicture}[
			node distance=1.1cm,
			neuron/.style={circle, draw=bitnetColor, thick, minimum size=0.7cm},
			lyr/.style={font=\scriptsize\bfseries, text=neutralColor},
			]
			\node[neuron] (x1) at (0,1.2) {};
			\node[neuron] (x2) at (0,0) {};
			\node[neuron] (x3) at (0,-1.2) {};
			\node[neuron] (a) at (2.7,0.8) {};
			\node[neuron] (b) at (2.7,-0.8) {};
			\node[neuron] (c) at (5.4,0.8) {};
			\node[neuron] (d) at (5.4,-0.8) {};
			\node[neuron] (o) at (8.1,0) {};
			\foreach \i in {x1,x2,x3} \foreach \j in {a,b} \draw[-{Latex[length=1.5mm]}, draw=neutralColor!70] (\i) -- (\j);
			\foreach \i in {a,b} \foreach \j in {c,d} \draw[-{Latex[length=1.5mm]}, draw=neutralColor!70] (\i) -- (\j);
			\foreach \i in {c,d} \draw[-{Latex[length=1.5mm]}, draw=neutralColor!70] (\i) -- (o);
			\node[lyr, above=0.3cm of x1] {Entrada (3)};
			\node[lyr, above=0.3cm of a] {Camada 1 (2)};
			\node[lyr, above=0.3cm of c] {Camada 2 (2)};
			\node[lyr, above=0.3cm of o] {Saída (1)};
		\end{tikzpicture}
	\end{center}
	\par\small Esta é exatamente a rede usada nas demonstrações desta seção: 3 entradas, duas camadas ocultas de 2 neurônios, 1 saída, todas com ativação sigmoide.
	% [SOFTWARE] Backprop -> Rede de 4 camadas | mostra o forward neuronio a
	% neuronio na mesma topologia 3-2-2-1 deste diagrama
	\abrirNoSoftware{Backprop \textrightarrow{} Rede de 4 camadas}{backprop.mlp}
\end{frame}

		\DemoSlide{Backprop \textrightarrow{} Rede de 4 camadas}{backprop.mlp}{Percorra o forward neurônio a neurônio, na mesma topologia 3-2-2-1 do diagrama.}
		\begin{frame}[shrink]
	\frametitle{O que é um peso, e como ele é calculado}
	\small
	\begin{itemize}
		\item Um \textbf{peso} é um número real associado a uma conexão entre dois neurônios --- é o parâmetro que a rede efetivamente \textbf{aprende}. Os dados de entrada não mudam; os pesos, sim.\newline
		
		\item Dentro de um neurônio, os pesos entram na \textbf{combinação linear}: cada entrada é multiplicada pelo seu peso, e os produtos são somados.
	\end{itemize}
	
	\begin{equation*}
		z = \sum_i w_i \, x_i \qquad y = \sigma(z)
	\end{equation*}
	
	\par Exemplo numérico real com o primeiro neurônio da rede do slide anterior:\newline
	\par Entrada $x = [1{,}0,\ 0{,}5,\ -0{,}5]$
	\par Pesos $w = [0{,}4,\ -0{,}3,\ 0{,}6]$.
	
	\begin{equation*}
		z = 0{,}4 \!\cdot\! 1{,}0 + (-0{,}3) \!\cdot\! 0{,}5 + 0{,}6 \!\cdot\! (-0{,}5) = -0{,}05
		\quad\Longrightarrow\quad
		y = \sigma(-0{,}05) \approx 0{,}4875
	\end{equation*}
	
	\par Cada um dos outros neurônios repete exatamente esse cálculo, com sua própria linha de pesos. Antes do treino, os pesos partem de algum valor inicial (aqui, fixos); o que o treino faz é ajustá-los, um pouco a cada passo, para que a saída da rede se aproxime do que se quer.
\end{frame}

		\begin{frame}[shrink]
	\frametitle{Como uma rede neural é treinada}
	\small
	\begin{itemize}
		\item \textbf{1. Forward}: calcula a saída $y$ a partir da entrada e dos pesos atuais.
		\item \textbf{2. Perda (loss)}: compara $y$ com o alvo desejado, num único número que cresce quanto maior o erro.
		\item \textbf{3. Backward}: a regra da cadeia calcula $\partial L/\partial w$ --- o quanto a perda mudaria se cada peso mudasse um pouco.
		\item \textbf{4. Atualização}: cada peso anda um pequeno passo no sentido \textbf{contrário} ao gradiente.
	\end{itemize}
	\begin{equation*}
		w \leftarrow w - \eta \cdot \frac{\partial L}{\partial w}
	\end{equation*}
	\par Exemplo numérico real (neurônio único, $x=2{,}0$, alvo $=0{,}9$, $w_0=-1{,}0$, taxa de aprendizado $\eta=3{,}0$):
	\begin{equation*}
		z = w_0 x = -2{,}0 \ \Rightarrow\ y = \sigma(z) \approx 0{,}1192 \ \Rightarrow\ L = \tfrac12(y-\text{alvo})^2 \approx 0{,}3048
	\end{equation*}
	\begin{equation*}
		\frac{\partial L}{\partial w} \approx -0{,}1640
		\quad\Longrightarrow\quad
		w_1 = -1{,}0 - 3{,}0 \cdot(-0{,}1640) \approx -0{,}508
	\end{equation*}
	\par Repetindo esses quatro passos, iteração após iteração, $y$ converge para o alvo. É exatamente esse ciclo --- forward, perda, backward, atualização --- que o resto desta palestra questiona: o que acontece quando os pesos só podem valer $\{-1,0,1\}$ (BitNets), ou quando a ativação vira um pulso discreto no tempo (redes de pulso)?
	% [SOFTWARE] Backprop -> Forward e backward classicos | passo a passo dos
	% 9 estagios da iteracao 1 reproduz exatamente estes numeros
	\abrirNoSoftware{Backprop \textrightarrow{} Forward e backward clássicos}{backprop.classic}
\end{frame}

		\DemoSlide{Backprop \textrightarrow{} Forward e backward clássicos}{backprop.classic}{Percorra os 9 estágios da iteração 1 --- os números batem exatamente com os do slide anterior.}

	\section{Estrutura da apresentação}
		\begin{frame}
	\frametitle{Estrutura da apresentação}
	\begin{itemize}
		\item Motivação: o gargalo comum às duas fronteiras.
		\item Taxonomia: dois eixos ortogonais.
		\item Parte I --- BitNets: quantização extrema de pesos.
		\item Parte II --- Redes de pulso: computação orientada a evento.
		\item Convergência: onde as duas se encontram.
		\item Desafios abertos e síntese.
	\end{itemize}
\end{frame}


	\setbeamercolor{title}{fg=chromeText,bg=accentColor!75!white}
	\setbeamercolor{frametitle}{fg=chromeText,bg=accentColor!75!white}
	\section{Motivação}
		\begin{frame}
	\frametitle{O gargalo}
	\only<1>{
		\framesubtitle{Pouca memória muitos dados}
		\begin{itemize}
			\item Redes convencionais (RNAs dedicadas a tarefas de larga escala) crescem em parâmetros muito mais rápido que o poder de processamento dos chips necessitando uma grande quantidade de hardware.\newline
			
			\item Um LLM de 7 bilhões de parâmetros em FP16 ocupa $7\times10^{9}\times 2\text{ bytes} \approx 14\text{ GB}$ \textbf{só para os pesos} --- antes do cache de atenção e das ativações.
			\item Considerando que uma GPU de consumo comum tem 8--24 GB. O modelo não cabe.
		\end{itemize}
	}
	\only<2>{
		\framesubtitle{A implicação}
		\begin{block}{Duas respostas independentes, dois eixos diferentes}
			\begin{itemize}
				\item \textbf{Reduzir quantos bits} descrevem cada peso/ativação $\rightarrow$ \textcolor{bitnetColor}{\textbf{BitNets}}.
				\item \textbf{Reduzir quando/quanto se computa}, processando só em resposta a eventos esparsos $\rightarrow$ \textcolor{snnColor}{\textbf{Redes de pulso}}.
			\end{itemize}
		\end{block}
	}
\end{frame}

		\begin{frame}
	\frametitle{Duas fronteiras}
	\small
	\begin{columns}
		\column{.48\linewidth}
		\begin{beamercolorbox}[sep=5pt,rounded=true,shadow=true]{titleBitnet}
			\textbf{Fronteira 1 --- BitNets}
		\end{beamercolorbox}
		\vspace{1pt}
		\par Eixo da \textbf{representação}: quantos bits descrevem cada peso/ativação.
		\begin{itemize}
			\item FP32 $\rightarrow$ FP16 $\rightarrow$ INT8 $\rightarrow$ ternário $\{-1,0,1\}$ $\rightarrow$ binário $\{-1,1\}$.
			\item A rede continua densa e síncrona: \textbf{todo} neurônio ainda computa em \textbf{toda} passada.
			\item O que muda é o \textbf{custo} de cada computação individual.
		\end{itemize}
		\column{.48\linewidth}
		\begin{beamercolorbox}[sep=5pt,rounded=true,shadow=true]{titleSnn}
			\textbf{Fronteira 2 --- Redes de pulso}
		\end{beamercolorbox}
		\vspace{1pt}
		\par Eixo do \textbf{tempo/atividade}: quando e quanto se computa.
		\begin{itemize}
			\item Ativação contínua e síncrona $\rightarrow$ pulso esparso, assíncrono, orientado a evento.
			\item Os pesos podem continuar em precisão plena.
			\item O que muda é \textbf{se} um neurônio computa, a cada instante.
		\end{itemize}
	\end{columns}
	\begin{block}{O alvo é o mesmo}
		Ambas atacam o custo (energia, memória, latência) da multiplicação densa em ponto flutuante --- mas por caminhos diferentes.
	\end{block}
\end{frame}


	\setbeamercolor{title}{fg=chromeText,bg=bitnetColor!92!white}
	\setbeamercolor{frametitle}{fg=chromeText,bg=bitnetColor!92!white}
	\section{Parte I --- BitNets}
		\begin{frame}
	\frametitle{Linha do tempo: da quantização discreta às LLMs de 1 bit}
	\centering
	\resizebox{.98\linewidth}{!}{%
	\begin{tikzpicture}[
		yr/.style={draw=bitnetColor, fill=bitnetColor!15, rounded corners=2pt, minimum height=0.9cm, minimum width=1.6cm, font=\tiny\sffamily, align=center},
		lb/.style={font=\tiny\sffamily, align=center, text width=1.7cm}
		]
		\draw[-{Latex[length=2mm]}, line width=0.6mm, draw=neutralColor] (-0.3,0) -- (12.3,0);

		\node[yr] (n1) at (0,0) {2015};
		\node[lb, below=2pt of n1] {BinaryConnect \cite{courbariaux2016binarized}};

		\node[yr] (n2) at (2.4,0) {2016};
		\node[lb, above=2pt of n2] {BNN \\ XNOR-Net \cite{rastegari2016xnor} \\ TWN \cite{li2016ternary}};

		\node[yr] (n3) at (4.8,0) {2017};
		\node[lb, below=2pt of n3] {TTQ \cite{zhu2017trained}};

		\node[yr] (n4) at (7.2,0) {2022};
		\node[lb, above=2pt of n4] {LLM.int8() \cite{dettmers2022llm}};

		\node[yr] (n5) at (9.6,0) {2023};
		\node[lb, below=2pt of n5] {GPTQ \cite{frantar2023gptq} \\ SmoothQuant \cite{xiao2023smoothquant} \\ BitNet \cite{wang2023bitnet}};

		\node[yr, fill=bitnetColor!45] (n6) at (12,0) {2024};
		\node[lb, above=2pt of n6] {\textbf{BitNet b1.58} \cite{ma2024era}};
	\end{tikzpicture}%
	}
	\vspace{2pt}
	\begin{itemize}\small
		\item 2015--2017: quantização extrema nasce em \textbf{visão computacional} (redes pequenas, treinadas do zero).
		\item 2022--2023: quantização \textbf{pós-treino} (PTQ) vira estratégia dominante para LLMs já treinados.
		\item 2023--2024: BitNet retoma a ideia de \textbf{treinar do zero} em baixa precisão --- agora em escala de bilhões de parâmetros.
	\end{itemize}
\end{frame}

		\begin{frame}
	\frametitle{Quantizar depois, ou quantizar desde o início?}
	\only<1>{
		\framesubtitle{Dois caminhos, confundidos com frequência}
		\begin{columns}
			\column{.48\linewidth}
			\begin{beamercolorbox}[sep=5pt,rounded=true,shadow=true]{title}
				\textbf{PTQ}: Quantização pós-treino
			\end{beamercolorbox}
			\begin{itemize}
				\item Treina em FP32/FP16 normalmente.
				\item Converte os pesos \textbf{depois}, sem retreinar (ou com pouco ajuste).
				\item Rápido e barato: LLM.int8(), GPTQ, SmoothQuant.
				\item A rede nunca ``viu'' o erro de quantização durante o treino.
			\end{itemize}
			\column{.48\linewidth}
			\begin{beamercolorbox}[sep=5pt,rounded=true,shadow=true]{title}
				\textbf{QAT}: Quantização durante o treino
			\end{beamercolorbox}
			\begin{itemize}
				\item A quantização já ocorre no \textit{forward pass} do treino.
				\item BinaryConnect, BNN, TWN, TTQ, BitNet, BitNet b1.58.
				\item Caro: em geral exige treinar do zero.
				\item A rede se \textbf{adapta} ao erro de quantização.
			\end{itemize}
		\end{columns}
	}
	\only<2>{
		\framesubtitle{O modo de falha da PTQ agressiva}
		\small
		\begin{itemize}
			\setlength{\itemsep}{0pt}
			\item Abaixo de $\sim$4 bits, a acurácia de modelos quantizados por PTQ pode cair \textbf{abruptamente} --- não gradualmente.
			\item Mecanismo concreto \cite{xiao2023smoothquant}: algumas ativações de LLMs têm \textbf{outliers} --- valores até $100\times$ maiores que o resto do tensor, concentrados em poucos canais.
			\item Um único fator de escala por tensor precisa cobrir o outlier \textbf{e} os valores normais $\rightarrow$ os valores normais perdem toda a resolução.
			\item Esse é o motivo de existir uma cadeia inteira de métodos de PTQ (LLM.int8(), GPTQ, SmoothQuant) --- cada um mitigando esse efeito de um jeito diferente, em vez de resolvê-lo evitando quantização baixa.
		\end{itemize}
		\begin{alertblock}{Onde o BitNet muda de estratégia}
			Se a rede é treinada \textbf{desde o início} em baixa precisão, ela nunca desenvolve ativações que dependam dessa resolução --- o problema é evitado, não corrigido depois.
		\end{alertblock}
	}
\end{frame}

		\begin{frame}
	\frametitle{A camada BitLinear}
	\only<1>{
		\framesubtitle{Da camada Linear densa à ternária}
		\par Uma camada Linear comum calcula $Y = W X$, com $W$ em FP16/FP32: cada saída soma \textbf{multiplicações reais} entre pesos e ativações contínuos.\newline
		
		\par O BitNet b1.58 \cite{ma2024era} substitui $W$ por uma versão ternária $\widetilde{W} \in \{-1,0,1\}$, obtida pela \textbf{quantização \textit{absmax}}:
		
		\begin{equation}
			\gamma = \text{mean}(|W|), \qquad
			\widetilde{W} = \text{round}\big(\text{clip}(W/\gamma,\,-1,\,1)\big)
		\end{equation}
		\par $\gamma$ é um único número por tensor --- guardado à parte para reescalar a saída no fim, não para cada peso individualmente.
	}
	\only<2>{
		\framesubtitle{Exemplo numérico}
		\par Tome três pesos de uma linha de $W$: $w = [0{,}6,\ -0{,}3,\ 0{,}05]$.
		\begin{equation*}
			\gamma = \frac{0{,}6+0{,}3+0{,}05}{3} = \frac{0{,}95}{3} = 0{,}3167
		\end{equation*}
		% quebrado em duas linhas: os tres estagios numa linha so estouravam a
		% largura do slide (Overfull \hbox de 14,5pt)
		\begin{equation*}
			\begin{aligned}
				w/\gamma &= [1{,}895,\ -0{,}947,\ 0{,}158]
				\xrightarrow{\text{clip}(-1,1)} [1{,}0,\ -0{,}947,\ 0{,}158] \\
				&\xrightarrow{\text{round}} \boxed{[1,\ -1,\ 0]}
			\end{aligned}
		\end{equation*}
		\par O peso original $0{,}05$ (o menor em módulo) \textbf{vira zero} --- é descartado, não aproximado. Os outros dois viram $\pm 1$.
		\par As ativações $X$ recebem seu próprio tratamento, quantização \textbf{absmax} para INT8: $\widetilde{X} = \text{round}(X \cdot 127/\max(|X|))$, com \textit{clip} em $[-128,127]$.
		% [SOFTWARE] BitNet -> Quantização | slider w ~ 0 para reproduzir a
		% zona morta, depois 0,6 / -0,3 para os dois outros pesos do exemplo
		\abrirNoSoftware{BitNet \textrightarrow{} Quantização}{bitnet.quant}
	}
	\only<3>{
		\framesubtitle{O que essa troca faz}
		\begin{itemize}
			\item Com $\widetilde{W} \in \{-1,0,1\}$, a multiplicação $\widetilde{W}\cdot\widetilde{X}$ vira: \newline 
			\textbf{some} $\widetilde{X}$ (peso $+1$), \newline
			\textbf{subtraia} $\widetilde{X}$ (peso $-1$), ou \newline
			\textbf{ignore} (peso $0$).\newline
			
			\item Nenhuma multiplicação real acontece no lado dos pesos --- só somas/subtrações de inteiros, mais um único rescalonamento por $\gamma$ ao final.
		\end{itemize}
	}
\end{frame}

		\DemoSlide{BitNet \textrightarrow{} Quantizacao}{bitnet.quant}{Ajuste o slider w para perto de 0 e reproduza a zona morta; depois experimente 0,6 e -0,3, os dois pesos do exemplo.}
		\DemoSlide{BitNet \textrightarrow{} Forward}{bitnet.forward}{Siga os checkpoints "Quantizar w1"/"w2" até "Multiplicação 1"/"2" --- Q(w2)=0 anula x2.}
		\begin{frame}[shrink]
	\frametitle{Treinar através de uma função que não tem gradiente}
	\begin{itemize}
		\item $\text{round}(\cdot)$ e $\text{clip}(\cdot)$ são \textbf{funções em degrau}: a derivada é zero em quase todo ponto, e indefinida nos degraus.
		\item Retropropagação padrão não tem o que propagar através delas --- o gradiente morre.
	\end{itemize}
	\vspace{1pt}
	\begin{block}{A solução: \textit{Straight-Through Estimator} (STE)}
		\begin{itemize}
			\item \textbf{Forward}: usa os pesos quantizados $\widetilde{W}$ normalmente.
			\item \textbf{Backward}: finge que a quantização foi a função identidade --- deixa o gradiente passar direto para uma cópia dos pesos em precisão plena, mantida ``nos bastidores'' e só quantizada de novo a cada passada.
		\end{itemize}
	\end{block}
	\vspace{0pt}
	\begin{alertblock}{Ponte com a Parte II}
		É a \textbf{mesma} ideia por trás do \textit{gradiente substituto} em redes de pulso \cite{eshraghian2023training}: forward com a função em degrau real, backward com uma aproximação suave, diferenciável. Duas subáreas, independentemente, chegaram à mesma solução para o mesmo problema estrutural.
	\end{alertblock}
\end{frame}

		\DemoSlide{BitNet \textrightarrow{} Backward \textrightarrow{} STE}{bitnet.ste}{Avance até o checkpoint final: derivada real vs. constante 1 do STE.}
		\begin{frame}
	\frametitle{BitNet b1.58}
	\framesubtitle{Comparado a LLaMA em FP16, mesmo tamanho de modelo \cite{ma2024era}}
	\begin{center}
		\begin{tabular}{lcc}
			\toprule
			& \textbf{3B parâmetros} & \textbf{70B parâmetros} \\
			\midrule
			Memória de GPU & $3{,}55\times$ menor & $7{,}2\times$ menor \\
			Velocidade de decodificação & $2{,}71\times$ mais rápido & $4{,}1\times$ mais rápido \\
			\textit{Throughput} & --- & $8{,}9\times$ maior \\
			Perplexidade / tarefas finais & \multicolumn{2}{c}{equivalente ao FP16} \\
			\bottomrule
		\end{tabular}
	\end{center}
	\vspace{2pt}
	\begin{itemize}
		\item Os autores reportam paridade estatística com a linha de base FP16 a partir de $\sim$3B parâmetros.\newline
		\item \alert{Abaixo dessa escala, a paridade some}: em modelos pequenos, a capacidade perdida pela quantização ternária pesa mais do que o modelo consegue compensar.\newline
		\item De onde vem o nome ``1,58 bit'': $\log_2(3) \approx 1{,}58$ --- o número de bits necessário para representar 3 estados ($-1,0,1$) por peso.
	\end{itemize}
\end{frame}

		\begin{frame}
	\frametitle{Hardware e limitações}
	\begin{columns}
		\column{.5\linewidth}
		\begin{beamercolorbox}[sep=5pt,rounded=true,shadow=true]{title}
			A promessa de hardware
		\end{beamercolorbox}
		\begin{itemize}
			\item Sem pesos em ponto flutuante, não há necessidade de multiplicadores de ponto flutuante --- só somadores de inteiros.
			\item Abre espaço para ASICs dedicados, mais simples e mais eficientes por operação.
		\end{itemize}
		\column{.5\linewidth}
		\begin{beamercolorbox}[sep=5pt,rounded=true,shadow=true]{title}
			Os limites atuais
		\end{beamercolorbox}
		\begin{itemize}
			\item Kernels/\textit{runtimes} para pesos ternários ainda são imaturos em GPUs convencionais, otimizadas para INT8/FP16 denso.
			\item Exige treinar do zero: não há como converter um checkpoint FP16 já treinado com a mesma qualidade.
			\item Ganho de paridade só aparece a partir de $\sim$3B parâmetros.
		\end{itemize}
	\end{columns}
\end{frame}


	\setbeamercolor{title}{fg=chromeText,bg=snnColor!85!white}
	\setbeamercolor{frametitle}{fg=chromeText,bg=snnColor!85!white}
	\section{Parte II --- Redes de pulso}
		\begin{frame}
	\frametitle{Redes neurais de pulso (RNP)}
	\only<1>{
		\framesubtitle{Por que pulsos?}
		\begin{itemize}
			\item Uma RNA convencional ativa \textbf{todos} os neurônios em \textbf{toda} passada --- o custo energético não depende do conteúdo do dado, só do tamanho da rede.
			\item O neurônio biológico só dispara quando o potencial acumulado ultrapassa um limiar; fica quieto o resto do tempo.
			\item Maass (1997) chamou as RNPs de \textbf{terceira geração} de modelos de redes neurais \cite{maass1997networks} --- depois do perceptron (1ª geração) e das redes com ativação contínua/sigmoide (2ª geração).
		\end{itemize}
	}
	\only<2>{
		\framesubtitle{Uma ressalva importante}
		\par Uma RNP \textbf{não} é uma simulação um-para-um de neurônios biológicos. Ela aproxima certas capacidades computacionais de propriedades biológicas específicas --- não o neurônio inteiro.
		\par Modelos mais fiéis à biologia existem (Hodgkin--Huxley \cite{gerstner2014neuronal}, dendritos não-lineares), mas são caros demais para treinar em escala. As RNPs usadas em aprendizado de máquina adotam modelos simplificados, como o \textbf{LIF} (próximo slide).
	}
\end{frame}

		\begin{frame}
	\frametitle{RNA vs. RNP: o par que se confunde}
	\small
	\begin{columns}
		\column{.48\linewidth}
		\begin{beamercolorbox}[sep=5pt,rounded=true,shadow=true]{title}
			Rede Neural Artificial (RNA)
		\end{beamercolorbox}
		\begin{itemize}
			\item Ativação: valor \textbf{contínuo} por neurônio, por camada.
			\item Processamento \textbf{síncrono}: uma passada completa por amostra.
			\item Sem dimensão temporal intrínseca (a menos que seja recorrente).
			\item Custo computacional fixo, independente da atividade.
		\end{itemize}
		\column{.48\linewidth}
		\begin{beamercolorbox}[sep=5pt,rounded=true,shadow=true]{title}
			Rede Neural de Pulso (RNP)
		\end{beamercolorbox}
		\begin{itemize}
			\item Ativação: \textbf{pulso} (0 ou 1), esparso no tempo.
			\item Processamento \textbf{orientado a evento}: só computa quando há pulso.
			\item Dinâmica temporal \textbf{intrínseca}: potencial de membrana acumula ao longo do tempo.
			\item Informação codificada em \textbf{taxa} e/ou \textbf{tempo} de disparo \cite{kasabov2019time}.
		\end{itemize}
	\end{columns}
	\vspace{2pt}
	\begin{alertblock}{Consequência prática}
		Uma RNP processa naturalmente \textbf{fluxos} de dados (áudio, EEG, eventos de câmera) --- mas processar um único dado estático exige ``esticá-lo'' ao longo de vários passos de tempo, um custo que a RNA não tem.
	\end{alertblock}
\end{frame}

		\DemoSlide{SNN \textrightarrow{} Sinal e spikes}{snn.spikes}{Aperte Play uma vez: o cruzamento de nível vira um pulso discreto.}
		\DemoSlide{SNN \textrightarrow{} Codificação Poisson}{snn.poisson}{Mesmo sinal da demo anterior, agora com pulsos probabilísticos --- compare.}
		\begin{frame}
	\frametitle{O neurônio \textit{Leaky Integrate-and-Fire} (LIF)}
	\only<1>{
		\framesubtitle{O modelo RC}
		\begin{columns}
			\column{.48\linewidth}
			\begin{itemize}
				\item O LIF se comporta como um circuito Resistor-Capacitor (RC).
				\item $R$: resistência de vazamento; $C$: capacitância; $V_{mem}$: potencial acumulado; $\vartheta$: interruptor que descarrega o capacitor (dispara um pulso) ao atingir o limiar.
				\item Governado por $\tau\frac{dV_{mem}}{dt} = R\,I_{in}(t) - V_{mem}(t)$, com $\tau = R\cdot C$ \cite{eshraghian2023training}.
			\end{itemize}
			\column{.48\linewidth}
			\begin{center}
				\includegraphics[width=.85\linewidth]{../../00-thesis/monography/images/rcmodel}\par
				{\tiny Fonte: \cite{eshraghian2023training}}
			\end{center}
		\end{columns}
	}
	\only<2>{
		\framesubtitle{Exemplo numérico: decaimento sem entrada}
		\par Com $R=5\,\Omega$, $C=1\,F$ (logo $\tau = R\cdot C = 5$) e potencial inicial $V_{mem}(0)=1$, sem corrente de entrada:
		\begin{equation*}
			V_{mem}(t) = V_{mem}(0)\cdot e^{-t/\tau}
			\quad\Longrightarrow\quad
			V_{mem}(1) = e^{-1/5} \approx 0{,}8187
		\end{equation*}
		\par O potencial \textbf{vaza} exponencialmente na ausência de estímulo --- daí o nome \textit{leaky}. Quando um estímulo chega, o potencial sobe (mesma matemática, sinal oposto) até cruzar um limiar $V_{th}$, disparando um pulso e reiniciando.
	}
	\only<3>{
		\framesubtitle{Comportamento completo simulado}
		\begin{center}
			\includegraphics[width=.62\linewidth]{../../00-thesis/monography/images/membranePotentialFull}\par
			{\tiny Corrente de 0,5 mA injetada entre os passos 51 e 70. Fonte: elaborado pelo autor.}
		\end{center}
	}
\end{frame}

		\DemoSlide{SNN \textrightarrow{} LIF}{snn.lif}{Os parâmetros padrão (tau=5, R=5, V\_th=1) já reproduzem o exemplo numérico anterior; aperte Play para o comportamento completo.}
		\begin{frame}
	\frametitle{Esparsidade: onde mora a eficiência}
	\only<1>{
		\begin{columns}
			\column{.55\linewidth}
			\begin{center}
				\includegraphics[height=3.3cm]{../../00-thesis/monography/images/sparsity}\par
				{\tiny Eixo horizontal: tempo. Eixo vertical: índice do neurônio. Fonte: elaborado pelo autor.}
			\end{center}
			\column{.42\linewidth}
			\small
			\begin{itemize}
				\item Na maior parte do tempo, a \textbf{grande maioria} dos neurônios está silenciosa.
				\item Pulsos são representados como \textbf{uns} dispersamente distribuídos em uma sequência de \textbf{zeros} \cite{eshraghian2023training}.
				\item Um neurônio LIF bem parametrizado tolera ruído: age como um \textbf{filtro passa-baixa}, só disparando quando o sinal se acumula de forma sustentada.
			\end{itemize}
		\end{columns}
		% [SOFTWARE] SNN -> Codificacao Poisson (imagem) | t ~ 15/29,
		% aponte que cada pixel e um "neuronio" sorteando disparo
		\abrirNoSoftware{SNN -> Codificação Poisson (imagem)}{snn.poisson_image}
	}
	\only<2>{
		\begin{alertblock}{A promessa, com precisão}
			Esparsidade só vira economia de energia real em hardware que sabe \textbf{pular} o zero. Em GPU/TPU convencional, a esparsidade é invisível ao hardware --- ele computa a multiplicação por zero do mesmo jeito.
		\end{alertblock}
		\vspace{4pt}
		\par Volta-se a este ponto na próxima seção, sobre hardware neuromórfico.
	}
\end{frame}

		\DemoSlide{SNN \textrightarrow{} Codificação Poisson (imagem)}{snn.poisson_image}{Avance até t~15/29 e aponte: cada pixel é um neurônio sorteando disparo.}
		\begin{frame}
	\frametitle{Como treinar uma função que não tem gradiente?}
	\only<1>{
		\framesubtitle{O mesmo problema estrutural do BitNet}
		\par A função de disparo do LIF se aproxima de uma \textbf{função em degrau} (Heaviside): dispara ou não dispara, sem meio-termo. Sua derivada é zero em quase todo ponto --- a descida de gradiente padrão não tem o que propagar \cite{kasabov2019time}.
		\par Quatro famílias de resposta, historicamente:
		\begin{itemize}
			\item \textbf{STDP}: regra local --- se o neurônio pré-sináptico dispara \textbf{antes} do pós-sináptico, a conexão se fortalece; se depois, se enfraquece. Biologicamente plausível, não supervisionada.
			\item \textbf{Gradiente substituto + BPTT}: usa a função em degrau real no \textit{forward}, e uma aproximação suave e diferenciável só no \textit{backward} --- a mesma lógica do STE do BitNet.
			\item \textbf{Algoritmos evolutivos}: seleção dos hiperparâmetros mais aptos ao longo de gerações.
			\item \textbf{Reservatório / computação dinâmica}: treina só a camada de saída, mantendo pesos recorrentes fixos (\textit{echo state networks}, \textit{liquid state machines}).
		\end{itemize}
	}
	\only<2>{
		\framesubtitle{Alternativa: converter em vez de treinar}
		\begin{itemize}
			\item Treina uma \textbf{RNA convencional}, com técnicas padrão de otimização.
			\item \textbf{Converte} os pesos treinados para uma RNP, balanceando pesos e limiares de disparo \cite{diehl2015fast, rueckauer2017conversion}.
			\item Evita o problema do gradiente por completo --- mas herda limitações: geralmente precisa de \textbf{mais passos de tempo} para aproximar a acurácia da RNA original, e nem toda arquitetura converte bem.
		\end{itemize}
	}
	% [SOFTWARE] SNN -> Surrogate gradient | checkpoints "A derivada
	% real" -> "O gradiente substituto" (sweep sigmoide+gradiente)
	\abrirNoSoftware{SNN \textrightarrow{} Surrogate gradient}{snn.surrogate}
\end{frame}

		\DemoSlide{SNN \textrightarrow{} Surrogate gradient}{snn.surrogate}{Percorra os checkpoints "A derivada real" até "O gradiente substituto".}
		\begin{frame}[shrink]
	\frametitle{Retropropagação através do tempo (BPTT)}
	\small
	\begin{columns}[T]
		\column{.55\linewidth}
		\resizebox{!}{3.3cm}{				\begin{tikzpicture}[
	box/.style={draw, rectangle, rounded corners=3pt,
		minimum width=1.2cm, minimum height=1.0cm,
		fill=gray!15, font=\small},
	fwd/.style={->, thick, blue!65!black},
	bwd/.style={->, thick, dashed, red!65!black},
	lbl/.style={font=\footnotesize}
	]
	% Forma compacta
	\node[box] (rc) {$\mathbf{W}$};
	\coordinate (xc) at ([yshift=0.9cm]rc.north);
	\node[lbl, above] at (xc) {$x_t$};
	\draw[fwd] (xc) -- (rc.north);
	\coordinate (sc) at ([yshift=-0.9cm]rc.south);
	\node[lbl, below] at (sc) {$s_t$};
	\draw[fwd] (rc.south) -- (sc);
	\draw[fwd] (rc.north east)
	to[out=50,in=-50,looseness=3.5]
	node[lbl, right] {$h_t$}
	(rc.south east);
	\node[lbl, below=2.4cm of rc, align=center] {Forma\\compacta};
	
	% Seta de equivalência
	\node[right=1.6cm of rc, font=\Large] {$\Rightarrow$};
	
	% Forma desenrolada
	\node[box, right=3.0cm of rc] (r1) {$\mathbf{W}$};
	\node[box, right=2.2cm of r1]  (r2) {$\mathbf{W}$};
	\node[lbl, right=1.4cm of r2]  (rd) {$\cdots$};
	\node[box, right=1.4cm of rd]  (rT) {$\mathbf{W}$};
	
	% Entradas
	\foreach \nd/\xi in {r1/1, r2/2, rT/T}{
		\coordinate (inp\nd) at ([yshift=0.9cm]\nd.north);
		\node[lbl, above] at (inp\nd) {$x_{\xi}$};
		\draw[fwd] (inp\nd) -- (\nd.north);
	}
	
	% Estado recorrente (forward)
	\draw[fwd] (r1.east) -- node[lbl, above] {$h_1$} (r2.west);
	\draw[fwd] (r2.east) -- node[lbl, above] {$h_2$} (rd.west);
	\draw[fwd] (rd.east) -- node[lbl, above] {$h_{T-1}$} (rT.west);
	
	% Saídas e perdas
	\foreach \nd/\xi in {r1/1, r2/2, rT/T}{
		\node[lbl, below=0.8cm of \nd] (L\xi) {$\mathcal{L}_{\xi}$};
		\draw[fwd] (\nd.south) -- (L\xi);
	}
	
	% Retropropagação do gradiente
	\draw[bwd] ([yshift=-0.3cm]r2.west) -- ([yshift=-0.3cm]r1.east);
	\draw[bwd] ([yshift=-0.3cm]rd.west) -- ([yshift=-0.3cm]r2.east);
	\draw[bwd] ([yshift=-0.3cm]rT.west) -- ([yshift=-0.3cm]rd.east);
	
	\node[lbl, below=1.6cm of r2, text=red!65!black]
	{$\longleftarrow$ gradiente retropropagado};
	\node[lbl, below=2.3cm of r2]
	{pesos $W$ compartilhados em todas as cópias};
	
\end{tikzpicture}}
		\column{.4\linewidth}
		\begin{itemize}
			\item A rede é \textbf{desenrolada} em $T$ cópias no tempo, compartilhando os mesmos pesos $W$.
			\item \textit{Forward}: o estado $h_t$ avança da esquerda para a direita.
			\item \textit{Backward}: o gradiente percorre o caminho inverso, acumulando a contribuição de cada passo:
		\end{itemize}
		\begin{equation*}
			\frac{\partial \mathcal{L}}{\partial W} = \sum_{t=1}^{T} \frac{\partial \mathcal{L}_t}{\partial W}
		\end{equation*}
	\end{columns}
	\vspace{2pt}
	\par\sloppy É neste \textit{backward}, a cada passo $t$, que o gradiente substituto entra: ele troca a derivada (nula) da função de disparo real por uma aproximação suave, deixando o gradiente fluir de volta pelos $T$ passos \cite{wu2018spatio, eshraghian2023training}.
	% [SOFTWARE] SNN -> Surrogate gradient | mesma demo de snnTreinamento.tex
	\abrirNoSoftware{SNN \textrightarrow{} Surrogate gradient}{snn.surrogate}
\end{frame}

		\label{sec:snnHardware}
		\begin{frame}
	\frametitle{Hardware neuromórfico}
	\begin{itemize}
		\item \textbf{TrueNorth} (IBM, 2014): 1 milhão de neurônios, chip totalmente assíncrono e orientado a evento \cite{merolla2014truenorth}.
		\item \textbf{Loihi} (Intel, 2018): aprendizado \textit{on-chip}, arquitetura \textit{manycore} orientada a evento \cite{davies2018loihi}.
		\item \textbf{SpiNNaker} (Manchester, 2014): milhões de núcleos ARM emulando pulsos em tempo real biológico \cite{furber2014spinnaker}.
	\end{itemize}
	\vspace{8pt}
	\begin{alertblock}{O modo de falha silencioso}
		O ganho energético de uma RNP só se realiza \textbf{nesse hardware especializado}. Rodada em GPU/TPU convencional --- otimizada para multiplicação densa, não para pular zeros --- simular $T$ passos de tempo por amostra pode deixar a RNP \textbf{mais lenta} que a RNA equivalente.
	\end{alertblock}
\end{frame}

		\begin{frame}
	\frametitle{Aplicações e redes de pulso + \textit{transformers}}
	\begin{itemize}
		\item \textbf{Câmeras de eventos} (DVS): sensores que já emitem pulsos nativamente --- casamento direto com RNPs.
		\item \textbf{EEG/BCI}: sinais biológicos já são fluxos temporais esparsos por natureza --- domínio onde a codificação temporal da RNP é vantagem, não custo extra. Para o formalismo de modelos de pulso em neurociência, ver \cite{dan_goodman_2022_7044500}.
		\item \textbf{Computação de borda}: dispositivos vestíveis e embarcados, onde bateria é o recurso mais escasso.
		\item \textbf{Spikformer} \cite{zhou2023spikformer}: reformula o mecanismo de atenção do \textit{transformer} em operações compatíveis com pulsos, aproximando as duas linhagens de arquitetura que dominam \textit{deep learning} hoje.
	\end{itemize}
	\vspace{8pt}
	\par Panorama mais amplo do campo em \cite{roy2019towards}.
\end{frame}


	\setbeamercolor{title}{fg=chromeText,bg=convergeColor!80!white}
	\setbeamercolor{frametitle}{fg=chromeText,bg=convergeColor!80!white}
	\section{Convergência}
		\label{sec:convergencia}
		\begin{frame}
	\frametitle{Onde as duas fronteiras se encontram}
	\centering
	\resizebox{!}{4.3cm}{%
	\begin{tikzpicture}[
		quad/.style={draw=neutralColor, minimum width=4.6cm, minimum height=3cm, align=center, font=\small},
		axislbl/.style={font=\small\sffamily, align=center}
		]
		% Quadrant fills
		\node[quad, fill=neutralColor!12] (bl) at (0,0) {\textbf{RNA padrão}\\[2pt]\footnotesize FP32/FP16, densa,\\ síncrona};
		\node[quad, fill=bitnetColor!18] (br) at (4.8,0) {\textbf{BitNet / BNN / TWN}\\[2pt]\footnotesize binário/ternário, mas\\ ainda densa e síncrona};
		\node[quad, fill=snnColor!18] (tl) at (0,3.2) {\textbf{RNP de precisão plena}\\[2pt]\footnotesize pesos FP32,\\ esparsa e orientada a evento};
		\node[quad, fill=convergeColor!25] (tr) at (4.8,3.2) {\textbf{RNP quantizada}\\[2pt]\footnotesize binário/ternário \emph{e}\\ esparsa/orientada a evento};

		% Axes
		\draw[-{Latex[length=2mm]}, draw=neutralColor!70] (-2.7,-1.6) -- (7.5,-1.6);
		\node[axislbl] at (2.4,-2.05) {precisão dos pesos: alta $\longrightarrow$ ternária/binária};
		\draw[-{Latex[length=2mm]}, draw=neutralColor!70] (-2.7,-1.6) -- (-2.7,4.8);
		\node[axislbl, rotate=90] at (-3.15,1.6) {denso/síncrono $\longrightarrow$ esparso/evento};
	\end{tikzpicture}%
	}
	\vspace{2pt}
	\begin{itemize}\small
		\item Um pulso já \textbf{é}, por definição, um valor binário no tempo (disparou ou não) --- a RNP é quantização extrema no eixo \textbf{temporal}, assim como o BitNet é quantização extrema no eixo dos \textbf{pesos}.
		\item O quadrante superior direito --- RNPs com pesos binários/ternários --- é direção de pesquisa \textbf{emergente}, ainda sem método canônico consolidado.
	\end{itemize}
	% [SOFTWARE] Comparacao -> ANN x BitNet x SNN | ultimo checkpoint
	% (tabela completa)
	\abrirNoSoftware{Comparação -> ANN x BitNet x SNN}{comparison}
\end{frame}

		\DemoSlide{Comparação \textrightarrow{} ANN x BitNet x SNN}{comparison}{Avance até o último checkpoint: a tabela completa ANN x BitNet x SNN.}
		\begin{frame}
	\frametitle{O tema comum}
	\begin{itemize}
		\item Ambas as fronteiras abandonam a \textbf{multiplicação densa em ponto flutuante} como operação fundamental --- uma pela via da \textbf{representação} (menos bits por peso), outra pela via do \textbf{tempo} (menos ativações por passo).
		\item As duas nasceram de comunidades diferentes (compressão de modelos vs. neurociência computacional) e, de forma independente, chegaram à mesma solução estrutural para treinar através de uma função não-diferenciável: o STE do BitNet e o gradiente substituto da RNP são o mesmo truque.
		\item O mesmo gargalo de fundo motiva as duas: energia dominada por \textbf{movimento de dado}, não por operações aritméticas em si --- a tabela de energia vista na abertura desta palestra.
	\end{itemize}
\end{frame}


	\setbeamercolor{title}{fg=chromeText,bg=accentColor!75!white}
	\setbeamercolor{frametitle}{fg=chromeText,bg=accentColor!75!white}
	\section{Desafios e síntese}
		\begin{frame}
	\frametitle{Desafios abertos}
	\begin{columns}
		\column{.48\linewidth}
		\begin{beamercolorbox}[sep=5pt,rounded=true,shadow=true]{title}
			\textcolor{bitnetColor}{\textbf{BitNets}}
		\end{beamercolorbox}
		\begin{itemize}
			\item Ecossistema de \textit{kernels}/hardware ainda imaturo.
			\item Exige treinar do zero --- não reaproveita checkpoints existentes.
			\item Paridade com FP16 só a partir de $\sim$3B parâmetros.
		\end{itemize}
		\column{.48\linewidth}
		\begin{beamercolorbox}[sep=5pt,rounded=true,shadow=true]{title}
			\textcolor{snnColor}{\textbf{Redes de pulso}}
		\end{beamercolorbox}
		\begin{itemize}
			\item \textit{Gap} de acurácia em tarefas complexas frente a RNAs grandes.
			\item Sobrecarga de simular $T$ passos em hardware convencional.
			\item Fragmentação entre \textit{frameworks} (snnTorch, SpikingJelly, Lava).
		\end{itemize}
	\end{columns}
	\vspace{10pt}
	\begin{alertblock}{Convergência}
		Co-projeto algoritmo-hardware ainda é incipiente para as duas frentes \textbf{separadamente} --- combiná-las (RNP quantizada) multiplica essa dificuldade em vez de somá-la.
	\end{alertblock}
\end{frame}

		\begin{frame}
	\frametitle{Síntese}
	\begin{itemize}
		\item Duas respostas independentes à mesma pressão: o custo energético e de memória do paradigma denso em ponto flutuante.
		\item \textcolor{bitnetColor}{\textbf{BitNets}} atacam a \textbf{representação} --- quantos bits descrevem cada peso.
		\item \textcolor{snnColor}{\textbf{Redes de pulso}} atacam o \textbf{tempo/atividade} --- quando vale a pena computar.
		\item O ganho real de ambas depende de hardware co-desenhado para explorá-las --- sem isso, permanecem truques de compressão em vez de mudanças de paradigma.
	\end{itemize}
\end{frame}

\begin{frame}
	\centering
	\Huge Obrigado.
	\vspace{1cm}

	\normalsize André Furlan --- ensismoebius@gmail.com
\end{frame}


	\begin{frame}[allowframebreaks]
		\frametitle{Referências}
		\bibliography{referencias.bib}
	\end{frame}

\end{document}
